\subsubsection{需求结构与机型容量的对应}
逐箱聚合如表\ref{tab:rev_cargo}。不同物资的质量与体积占比并不相同：饮用水构成主要质量负担，生活卫生用品则具有较高的单位质量体积。医疗物资的质量较小，但其硬截止对任务顺序具有直接影响。因此本文独立保留质量、体积和时限三个维度，不用平均货物密度或平均期限合并。
\begin{table}[htbp]\centering\small
\caption{逐箱需求的类型结构与容量规模}\label{tab:rev_cargo}
\begin{tabular}{lrrrr}\toprule
物资类型 & 箱数 & 质量/kg & 体积/m$^3$ & 质量占比/\%\\\midrule
医疗物资 & 16 & 48 & 0.192 & 6.33\\
应急食品 & 19 & 152 & 0.532 & 20.05\\
饮用水 & 36 & 504 & 0.972 & 66.49\\
生活卫生用品 & 9 & 54 & 0.315 & 7.12\\
\bottomrule\end{tabular}
\end{table}


为识别一个候选箱组首先触发的限制，定义容量占用比
\begin{equation}
 \rho_{rm}^{Q}=\frac{\sum_{b\in\mathcal B_r}\mu_b}{Q_m},\quad
 \rho_{rm}^{V}=\frac{\sum_{b\in\mathcal B_r}\nu_b}{V_m},\quad
 \rho_{rm}^{E}=\frac{E_r}{(1-\alpha_m)E_m}.
 \label{eq:dimensionless_capacity}
\end{equation}
三个比值均不大于1才满足当前任务的相应约束。$\rho^V$较大时应优先考虑货舱匹配；$\rho^E$较大时，瓶颈来自航段和载荷能耗；$\rho^Q$较大则反映额定质量接近上限。占用比用于解释约束，不替代原始单位下的可行性计算。

数据关联分别核对编号、量纲与总量守恒。每个箱号只能对应一个目的地和一套原始属性，所用目的地与机型必须存在于输入表中；组批及路线调整后，原始箱号全集、758 kg质量与2.011 m$^3$体积保持不变。优化改变的是任务组合与执行时序，不改变需求记录本身。
